\section{Introduction}\label{sec:introduction}
Fix \(0<\rho<1\), an odd integer \(L\ge5\), \(0<\alpha<1\), and
\(\varepsilon>0\). At level \(k\ge0\), put \(\eta=\eta_k=L^{-k}\).
A photon site is \(c\in\eta\Z^3/\ell\Z^3\); a chamber site is
\(x=(c,g)\), where \(g\in O_{\mathrm{cub}}\) is a signed permutation
matrix. There are \(48\) labels \(g\) and four spin components at
each photon site. Connections are periodic complex oriented bond
functions, with \(A(c',c)=-A(c,c')\); spinors are antiperiodic in
each coordinate. The chamber measure is \(\eta^3/48\), while the
photon-site spinor measure is \(\eta^3\).
Operator kernels use the chamber measure in their column variable;
the identity kernel is \(48\eta^{-3}\mathbf1[x=x']1_4\).
The cluster injection and mean are
\[
 (J\xi)(c,g)=\xi(c),\qquad J^*f(c)=48^{-1}\sum_g f(c,g).
\]
Adjoints use these Hermitian pairings. Transposes use the corresponding
bilinear pairings, without complex conjugation.

Use the Pauli and Dirac matrices
\[
 \sigma_1=\begin{pmatrix}0&1\\1&0\end{pmatrix},\quad
 \sigma_2=\begin{pmatrix}0&-i\\i&0\end{pmatrix},\quad
 \sigma_3=\begin{pmatrix}1&0\\0&-1\end{pmatrix},\qquad
 \gamma_j=\sigma_3\otimes\sigma_j,\quad
 \gamma_5=\sigma_1\otimes1_2.
\]
Let \(s_0\) reverse the third coordinate, \(s_1\) interchange
coordinates \(2,3\), and \(s_2\) interchange \(1,2\); let \(r_{e_1}\)
reverse the first coordinate. The four neighbours, inward normals,
and weights are
\[
 \begin{gathered}
 x^{(i)}=(c,gs_i)\quad(0\le i\le2),\qquad
 x^{(3)}=(c+\eta ge_1,gr_{e_1}),\\
 (n_0,n_1,n_2,n_3)
   =(e_3,(e_2-e_3)/\sqrt2,(e_1-e_2)/\sqrt2,-e_1),\\
 (w_0,w_1,w_2,w_3)=(6,6\sqrt2,6\sqrt2,6).
 \end{gathered}
\]
For a charge \(e_k>0\), the chamber operator is
\begin{equation}\label{eq:intro-chamber-operator}
 (D_Af)(x)=-\eta^{-1}\sum_{i=0}^3
       \frac{w_i}{2}(\rho+\gamma\cdot gn_i)
       \bigl(U_A(x,x^{(i)})f(x^{(i)})-f(x)\bigr),
\end{equation}
where internal transports are one and the external transport is
\(\exp(ie_k\eta A(c,c+\eta ge_1))\).

The unit blocks are
\(\Delta_z=\{c:|c-z|_\infty<1/2\}\), \(z\in\Z^3/\ell\Z^3\).
For \(c\in\Delta_z\), let \(c=c_0,c_1,\ldots,c_k=z\) be the
centres of its nested blocks of sides \(L^j\eta\).
Write \(\tau A(u,v)\) for the arithmetic mean of the oriented
fine-bond sums along the six coordinate-order paths from \(u\) to \(v\).
Set
\[
 \begin{aligned}
 \varphi_A(z,c)&=\exp\!\left(ie_k\eta
                    \sum_{j<k}\tau A(c_{j+1},c_j)\right),&
 (Q_k(A)\xi)(z)&=\eta^3\sum_{c\in\Delta_z}\varphi_A(z,c)\xi(c),\\
 (Q_k^T(-A)v)(c)&=\varphi_A(z,c)^{-1}v(z).&&
 \end{aligned}
\]
For \(z'\in L\Z^3/\ell\Z^3\), put \(B(z')=\{z:|z-z'|_\infty<L/2\}\).
The next average and its algebraic reverse are
\[
 \begin{aligned}
 U_A(z',z)&=\exp(ie_k\eta\tau A(z',z)),\\
 (Q(A)v)(z')&=L^{-3}\sum_{z\in B(z')}U_A(z',z)v(z),&
 (Q^T(-A)v')(z)&=U_A(z',z)^{-1}v'(z').
 \end{aligned}
\]
The target pairing has measure \(L^3\). These maps act componentwise
in spin. For \(k\ge1\), define
\[
 \Pi_1(A)=JQ_k^T(-A)Q_k(A)J^*,\qquad
 \Pi_L(A)=JQ_k^T(-A)Q^T(-A)Q(A)Q_k(A)J^*.
\]
At \(k=0\) use only \(\Pi_L(A)=JQ^T(-A)Q(A)J^*\).
Let \(b=1\), \(b_L=L^{-1}\), \(b_-=1-L^{-1}\), and
\[
 b_k=\frac{1-L^{-1}}{1-L^{-k}}\quad(k\ge1),\qquad
 \alpha_y=\frac{b_ki\gamma_5y}{b_k+i\gamma_5y},\qquad
 \beta_y=\frac{b_k^2b_L}{(b_k+b_L+i\gamma_5y)(b_k+i\gamma_5y)}.
\]
At \(y=\infty\), put \(\alpha_\infty=b_k\), \(\beta_\infty=0\).
For \(0<\mu\le1\), the penalized operators are
\begin{equation}\label{eq:intro-penalized-operator}
 F_y(A)=
 \begin{cases}
 D_A+\mu+\alpha_y\Pi_1(A)+\beta_y\Pi_L(A),
      & k\ge1,\ y\in\R\cup\{\infty\},\\
 D_A+\mu+b_L\Pi_L(A)+i\gamma_5y,
      & k=0,\ y\in\R.
 \end{cases}
\end{equation}
The data \(k,\mu,\ell,e_k\) are independent.

Let \(M\ge L^2\) be a power of \(L\), with partition \(M\)-cubes
centred in \(M\Z^3/\ell\Z^3\), and assume \(\ell/(LM)\) is odd
and at least \(13\). A connection belongs to
\(\widetilde{\mathcal G}_k\) if \(A=A_0+A_1\), \(A_0\) is real,
\(\|A_1\|_\infty<e_k^{-1+\varepsilon}\), and on every concentric
\(11M\) enlargement of a partition cube,
\[
 A_0=B_0+d\lambda,\qquad B_0,\lambda\ \hbox{real},\qquad
 \|B_0\|_\infty<e_k^{-1+\varepsilon},\qquad
 d\lambda(c,c')=\eta^{-1}(\lambda(c')-\lambda(c)).
\]
The local gauges may be chosen independently.

Distances are periodic sup distances. For distinct chamber sites set
\(d_X(x,x')=\max(d(c,c'),\eta)\), and, on pairs with \(0<d_X<1\), put
\[
 U_A(x,x')=\exp(ie_k\eta\tau A(c,c')),\qquad
 (\delta_{\alpha,A}v)(x,x')
 =\frac{U_A(x,x')v(x')-v(x)}
        {d_X(x,x')^\alpha}.
\]
Thus distinct labels over one photon site are included when \(k\ge1\).
For an operator \(T\), define
\[
 N_\gamma(T)=\sup_z\sum_{z'}e^{\gamma d(z,z')}
       \|\mathbf1_{\Delta_z}T\mathbf1_{\Delta_{z'}}\|_{\infty\to\infty}.
\]
The sup norm uses Euclidean spinor norms; for pair outputs the left
indicator restricts the first point. At level zero,
\(N_\gamma^{\mathrm{ker}}\) replaces the localized operator norm by
the supremum of the normalized chamber kernel on the two clusters.

\begin{theorem}\label{thm:chamber-locality}
There are \(a'>0\) and \(0<\gamma\le(6L)^{-1}\), depending only on
\(\rho,L,\alpha\), with the following property. Set
\(C_V=2\exp(1)(42+6b)\) and
\(q(p)=54\pi C_Va'p/M\). Assume
\begin{equation}\label{eq:main-locality-conditions}
 M>4L,\qquad q(1)\le\tfrac12,\qquad
 0<e_k\le e_0,\qquad
 6Le_0^\varepsilon\le1,\qquad 4C_Va'e_0^\varepsilon\le1.
\end{equation}
For every \(A\in\widetilde{\mathcal G}_k\), all the operators
\eqref{eq:intro-penalized-operator} are invertible. Their inverses
\(S_y(A)\) are analytic in the connection and satisfy
\[
 N_\gamma(S_y(A))\le32a'w_k(y),\qquad
 N_\gamma(\delta_{\alpha,A}S_y(A))\le304a',\qquad
 w_k(y)=
 \begin{cases}1,&k\ge1,\\(1+|y|)^{-1},&k=0.\end{cases}
\]
At level zero the first bound also holds in
\(N_\gamma^{\mathrm{ker}}\); there are no short pairs at that level.
The constants are independent of the level, mesh, volume,
permitted connection, reference mass \(\mu\in(0,1]\), and resolvent
parameter in its stated range.
\end{theorem}

We develop the chamber block arguments of
\cite[Sections~9 and~11]{QED3Manuscript}. The internal chamber graph
has a positive gap on fields with zero cluster mean. Eliminating those
components gives a four-component Schur symbol with the Dirac
singularity at small momentum. Its Fourier bounds supply local
regularity uniformly as the mesh tends to zero. Averaging penalties
control the remaining block means. Local real gauges then reduce the
connection to a small perturbation, and two local inverse identities
give convergent right and left walk expansions.

The walk supports also localize the coarse covariance and determinant.
We first attach parameters to the entire support of each walk, so
that setting parameters to zero separates components. Exact resolvent
identities pass from the fine inverse to the coarse inverse.
Weighted Berezin integration then gives the block formula in
\cref{thm:fermion-block-integration}. A resolvent integral expresses
its determinant as a sum of local analytic terms. For the additional
small-field difference estimate, long polymers supply the gain by
decay; on short polymers a real gauge gives an enlarged complex disk,
and Cauchy's formula gives the gain. The necessary nonwrapping
geometry is proved in \cref{app:polymer-geometry}.

The separate free estimates below hold for arbitrary unitary period
characters. The integration and transpose identities use the
antiperiodic convention fixed here. Their spin matrices are
\[
 \mathsf C=i\gamma_2,\qquad
 \varkappa=\gamma_5\mathsf C=\sigma_2\otimes\sigma_2,
 \qquad
 \mathsf C^T=\mathsf C^{-1}=-\mathsf C,\quad
 \varkappa^T=\varkappa=\varkappa^{-1}.
\]
Both matrices are real. Their displayed forms give the transpose and
complex-conjugation identities used in the determinant calculation.
The connection is an input to these estimates. For a connection
obtained from Maxwell reconstruction, the local gauge and fluctuation
domain estimates are supplied separately by
\cite[Theorems~A and~B]{MaxwellBlockManuscript}.
